%  Capítulo 0 de ejemplo — Este archivo NO se compila solo:
%  lo incluye main.tex con \chapter{Introducción}
En este documento se describe la aplicación Android que hemos desarrollado en Kotlin.
La aplicación reúne varias funcionalidades independientes a las que se accede desde un menú inicial: una calculadora, un conversor de divisas, una lista de recetas y un menú con tres juegos.

Los juegos son piedra, papel o tijera, un lanzador de dados y luces fuera.

Cada capítulo explica una funcionalidad, con una captura de la pantalla y las clases y funciones que hemos creado para ella.
En los casos más importantes se incluye también el fragmento de código.

Todas las pantallas comparten la misma paleta de colores para que la aplicación tenga un aspecto uniforme.

Al final del documento se incluye la información de uso de IA.

% chapter es un h1, section es un h2, subsection es un h3 y subsubsection es un h4.
% aqui puedes poner el titulo ñentre los corchetes y luego se pondra con el estilo en el pdf.
\chapter{Titulo 1}
% El label sirve para referenciarlo desde otro sitio con \ref{chap:intro}. sirve como un  enlace ñen lel pdf.                
\label{chap:ejemplo}

\section{Titulo 2}                        
Hola maya!.

Te pong aqui texto de ejemplo de todos los tipos para que puedas copiar y pegar.
Cada capítulo se pone en un archivo aparte y se incluye en el main.tex con input{capitulos/01-introduccion}.

Todo esto tienes que ediltarlo desde vscode, no desde el pdf.
Ahi veras estructuras de titulos, listas, etc.

Lo mejor para que no haya problemas es poner cada frase en una linea separada.
Un salto de línea simple no separa párrafos.
Para empezar un párrafo nuevo, deja una línea en blanco. (dos veces al intro)

\subsection{Titulo 3}                 
Puedes poner texto en \textbf{negrita}, en \textit{cursiva} o en \texttt{monospace}.

\subsubsection{Titulo 4}       
Este es el cuarto nivel de título.

\section{Listas}
% con esoto puedes hacer listas con viñetas
\begin{itemize}
  \item Lista con viñetas.
  \item Otro elemento.
% puedes meter una lista dentro de otra y cambia el estilo de la viñeta
  \begin{itemize}
    \item Elemento anidado.
  \end{itemize}
\end{itemize}

% con esto puedes hacer listas numeradas
\begin{enumerate}
  \item Lista numerada.
  \item Segundo paso.
\end{enumerate}


\section{Código}
Código en línea: \texttt{print("hola")}.

Bloque de código:
\begin{lstlisting}
def saludar(nombre):
    print(f"Hola, {nombre}")  # también con tildes: canción

saludar("mundo")
\end{lstlisting}

\section{Imágenes}
La imagen está en \texttt{figuras/ejemplo.png}; no hace falta poner la extensión.

\begin{figure}[H]
  \centering  % esto solo lo centra
  \includegraphics[width=0.6\textwidth]{ejemplo} % aqui se selecciona tanto la imagen como el tamaño. 0.6 es el 60% del ancho de la página.
  \caption{Descripción de la imagen.} % texto que aparece debajo de la imagen junto a "Figura 1: "
\end{figure}

Y también puedes referenciar capítulos: ver el capítulo~\ref{chap:intro}.
